% ICML 2026 poster -- De-Linearizing Agent Traces (BPOP)
% Built on the Gemini beamerposter theme: https://github.com/anishathalye/gemini

\documentclass[final]{beamer}

% ====================
% Packages
% ====================

\usepackage[T1]{fontenc}
\usepackage{lmodern}
% ICML 2026 main-conference suggested size: 36in (H) x 48in (W) = 91 x 122 cm (landscape)
\usepackage[size=custom,width=121.9,height=91.4,scale=1.2]{beamerposter}
\usetheme{gemini}
\usecolortheme{oxford}
\usepackage{graphicx}
\usepackage{booktabs}
\usepackage{anyfontsize}
\usepackage{amsmath, amssymb}
\usepackage{microtype}
% discourage very short last lines (widows) in justified paragraphs
\clubpenalty=10000 \widowpenalty=10000 \finalhyphendemerits=10000

% ====================
% Lengths
% ====================

% N=3 columns: (N+1)*sepwidth + N*colwidth = paperwidth
\newlength{\sepwidth}
\newlength{\colwidth}
\setlength{\sepwidth}{0.025\paperwidth}
\setlength{\colwidth}{0.3\paperwidth}

\newcommand{\separatorcolumn}{\begin{column}{\sepwidth}\end{column}}

% ====================
% Title
% ====================

\title{De-Linearizing Agent Traces with Bayesian Partial Orders\\[0.5ex]
  {\normalfont\LARGE Bayesian Inference of Latent Partial Orders for Efficient Execution}}

\author{Dongqing Li\textsuperscript{1,$\ast$} \and Zheqiao Cheng\textsuperscript{2,$\ast$}
        \and Geoff K.\ Nicholls\textsuperscript{1} \and Quyu Kong\textsuperscript{3}}

\institute[shortinst]{\textsuperscript{1}University of Oxford \quad
  \textsuperscript{2}Zhejiang University \quad
  \textsuperscript{3}Independent Researcher \quad
  ($\ast$\,equal contribution)}

% ====================
% Footer
% ====================

\footercontent{
  \href{https://arxiv.org/abs/2602.02806}{arXiv:2602.02806} \hfill
  ICML 2026, Seoul, South Korea \hfill
  \href{mailto:dongqing.li@kellogg.ox.ac.uk}{dongqing.li@kellogg.ox.ac.uk}}

% ====================
% Logos
% ====================
% Drop your files into figures/ and they are picked up automatically.
% Left  = Oxford : figures/oxplogo.png (or oxford-stats-logo.{pdf,png})
% Right = ICML 2026 : figures/icml-2026-logo.{pdf,png}
\IfFileExists{figures/oxplogo.png}
  {\logoleft{\hspace{2.5cm}\includegraphics[height=7cm]{figures/oxplogo.png}\hspace{0.8cm}%
     \raisebox{1.8cm}{\parbox[b]{15cm}{\fontsize{40}{46}\selectfont\bfseries Department of\\[0.2ex]Statistics}}}}
  {\IfFileExists{figures/oxford-stats-logo.pdf}
    {\logoleft{\includegraphics[height=7cm]{figures/oxford-stats-logo.pdf}}}
    {\IfFileExists{figures/oxford-stats-logo.png}
      {\logoleft{\includegraphics[height=7cm]{figures/oxford-stats-logo.png}}}{}}}
\IfFileExists{figures/icml-2026-logo.pdf}
  {\logoright{\includegraphics[height=7cm]{figures/icml-2026-logo.pdf}}}
  {\IfFileExists{figures/icml-2026-logo.png}
    {\logoright{\includegraphics[height=7cm]{figures/icml-2026-logo.png}}}{}}

% ====================
% Body
% ====================

\begin{document}

\begin{frame}[t]
\begin{columns}[t]
\separatorcolumn

% ============================================================
% COLUMN 1 -- Motivation & Idea
% ============================================================
\begin{column}{\colwidth}

  \begin{exampleblock}{The Big Idea}

    {\large Successful agent traces are \emph{not} recipes --- they are
    \textbf{noisy linearizations of a latent partial-order workflow}.}

    \begin{figure}
      \centering
      \includegraphics[width=0.40\colwidth]{figures/fig2_chaos.png}
      \caption{One latent plan (left) explains many valid execution orders
      (right); the \textbf{frontier} is the set of ready actions.}
    \end{figure}

  \end{exampleblock}

  \begin{block}{The Problem: Agents Re-Reason a Fixed Recipe}

    LLM agents execute workflows as \textbf{sequential traces}, which:
    \begin{itemize}
      \item \textbf{Hide concurrency} --- independent steps look ordered,
        blocking parallelism.
      \item \textbf{Trigger re-reasoning} every run (the ``perpetual intern'')
        --- slow and costly.
    \end{itemize}
    \textbf{BPOP:} recover the latent partial order \emph{once}, then compile it
    into a fast, reusable graph executor (the ``domain expert'').

  \end{block}

  \begin{block}{Traditional vs.\ BPOP Agent}

    \begin{figure}
      \centering
      \includegraphics[width=0.92\colwidth]{figures/architecture.pdf}
      \caption{BPOP recovers the partial order and compiles a low-latency,
      parallel graph executor.}
    \end{figure}

    \heading{Contributions}
    \begin{enumerate}
      \item Agent traces as \textbf{noisy linearizations} of a latent
        partial-order workflow.
      \item A tractable \textbf{frontier-softmax} likelihood that avoids
        \#P-hard marginalization over linear extensions.
      \item \textbf{Recoverability \& feasibility} analysis; open-sourced
        \textbf{Cloud-IaC-6} benchmark.
      \item \textbf{Graph compilation} into a frontier-based executor.
    \end{enumerate}

  \end{block}

\end{column}

\separatorcolumn

% ============================================================
% COLUMN 2 -- Method
% ============================================================
\begin{column}{\colwidth}

  \begin{block}{Model: Bayesian Partial Order Planning}

    \heading{Prior over partial orders}
    Each action $a_j$ gets a latent vector $U_j\in\mathbb{R}^K$. The order
    $h=(\mathcal{A},\succ_h)$ is the \textbf{intersection of $K$ total orders}:
    $$
    a_i \succ_h a_j \iff U_{i,k} > U_{j,k}\quad \forall k\in\{1,\dots,K\}.
    $$

    \begin{figure}
      \centering
      \includegraphics[width=0.44\colwidth]{figures/latent_U_crop.png}
      \caption{Latent coordinates $U_{j,k}$: Action 1 dominates all $k$, so it
      precedes every other action.}
    \end{figure}

    \heading{Frontier-softmax likelihood}
    With remaining actions $R_t$, the frontier is $F_t(h)=\{a\in R_t :
    \nexists\,b\in R_t,\ b\succ_h a\}$. A trace is a product of local choices:
    $$
    p(y_t\mid y_{<t},h)=(1-\epsilon)\,
      \frac{\exp\!\big(\beta\,Q(y_t)\big)}
           {\sum_{a\in F_t(h)}\exp\!\big(\beta\,Q(a)\big)}
      \;+\;\frac{\epsilon}{|R_t|},
    $$
    ($=0$ for $y_t\notin F_t$). $\beta$: inverse temperature; $\epsilon$:
    ``trembling-hand'' noise. Utility $Q(a)=\log(1+S_t(a))$ favors bottleneck
    actions ($S_t$ = descendant count).

    \textbf{Tractable:} $O(|\!\succ_h\!|+T|\mathcal{A}|)$ vs.\ \#P-hard counting
    of linear extensions.

    \begin{figure}
      \centering
      \includegraphics[width=0.95\colwidth]{figures/fig10_likelihood.png}
      \caption{Stepwise likelihood: at each $t$ the choice is scored over the
      frontier $F_t$ by utility $Q$; off-frontier picks ($y_3{=}4$) are explained
      only by the noise term $\epsilon/|R_t|$.}
    \end{figure}

  \end{block}

  \begin{block}{From Posterior to Executable Graph}

    MCMC over $(U,\beta,\epsilon)$ yields a posterior over $h$; threshold its
    relation probabilities at $\alpha$ to compile the SOP graph.
    \begin{itemize}
      \item $\alpha$ = \textbf{risk--efficiency knob}: higher $\Rightarrow$
        sparser, more parallel, riskier.
      \item Default $\alpha{=}1/3$ favors \textbf{recall} (safety-biased);
        $\widehat{h}_{\text{thr}}(1/3)\!\approx\!\widehat{h}_{\text{mode}}$.
      \item Compiled \textbf{Graph Execution Engine} fires ready actions in
        parallel, loading only relevant context per node.
    \end{itemize}

  \end{block}

\end{column}

\separatorcolumn

% ============================================================
% COLUMN 3 -- Results
% ============================================================
\begin{column}{\colwidth}

  \begin{block}{Structural Recovery vs.\ Trace Diversity}

    On \textbf{Cloud-IaC-6} ($n{=}54$ LLM traces) and WFCommons workflows vs.\
    Majority, Inductive/Heuristics Miner, Bayesian Queue-Jump.

    \begin{figure}
      \centering
      \includegraphics[width=0.66\colwidth]{figures/fig8_recovery.png}
      \caption{Edge-F1 vs.\ IP-Cov. BPOP improves with trace diversity and
      overtakes all baselines, while keeping high feasibility where trace-only
      methods are unstable or collapse to $0$.}
    \end{figure}

  \end{block}

  \begin{block}{Recovered SOP (Qualitative)}

    \begin{figure}
      \centering
      \includegraphics[width=0.80\colwidth]{figures/fig5_sop.png}
      \caption{Recovered vs.\ true SOPs (SRASearch, Epigenomics). Green: correct;
      red: missed (unsafe); orange: extra (safe). Errors are safety-biased.}
    \end{figure}

  \end{block}

  \begin{alertblock}{Key Results}

    \begin{itemize}
      \item \textbf{Best recovery:} Edge-F1 \textbf{0.91} (SRASearch) and
        \textbf{0.79} (Epigenomics), vs.\ $\le 0.43$ for every baseline.
      \item \textbf{Reasoning removed:} at IP-Cov $\ge 0.9$, completeness 100\%,
        fallback 0\%, cost $\to\!\sim$583 intent tokens/task
        (\textbf{0 reasoning tokens}) vs.\ pure-LLM's 234k--382k.
      \item \textbf{Hybrid} mode: only mode with 100\% success on all 6
        scenarios; errors are safety-biased (extra constraints, not missing).
    \end{itemize}

  \end{alertblock}

  \begin{block}{Paper, Project \& Posets}

    \begin{minipage}[t]{0.32\linewidth}
      \centering
      \includegraphics[width=0.44\linewidth]{figures/qr_paper.png}\\[0.3em]
      \textbf{Paper}\\ \footnotesize arXiv:2602.02806
    \end{minipage}\hfill
    \begin{minipage}[t]{0.32\linewidth}
      \centering
      \includegraphics[width=0.44\linewidth]{figures/qr_code.png}\\[0.3em]
      \textbf{Project \& Code}\\ \footnotesize github.com/hollyli-dq
    \end{minipage}\hfill
    \begin{minipage}[t]{0.32\linewidth}
      \centering
      \includegraphics[width=0.44\linewidth]{figures/project_qr.png}\\[0.3em]
      \textbf{Posets Project}\\ \footnotesize stats.ox.ac.uk/$\sim$nicholls
    \end{minipage}

  \end{block}

\end{column}

\separatorcolumn
\end{columns}
\end{frame}

\end{document}
